% ---------------------------------------------------------------------------
%  QML Lab research project 10: Quantum Reinforcement Learning on CartPole
%  Build:  latexmk -pdf main.tex   (needs qmlproject.sty, harvardqml.bst, references.bib)
% ---------------------------------------------------------------------------
\documentclass[11pt,a4paper]{article}
\usepackage{qmlproject}
\qmlheader{Project 10: Quantum RL on CartPole}

\hypersetup{pdftitle={QML Lab project 10: Quantum Reinforcement Learning on CartPole}}

\begin{document}

\begin{center}
  {\small\textcolor{black!55}{QML Lab \textbullet\ research project 10}}\\[0.5em]
  {\LARGE\bfseries Quantum Reinforcement Learning on CartPole}\\[0.6em]
  {\large A variational quantum circuit as the policy of a policy-gradient agent.}\\[0.6em]
  {\small Living document \textbullet\ last updated \today}
\end{center}

\begin{projectcard}
\small
\begin{tabularx}{\linewidth}{@{}>{\bfseries}l X >{\bfseries}l X@{}}
Quadrant  & \CQ\ (classical data, quantum processor) & Track & Hybrid \\
Difficulty & Hard & Tools & PennyLane, Gymnasium, PyTorch \\
Builds on & Modules 5 and 7 of the QML Lab & Workspace & \path{qml_projects/quantum-rl-cartpole/} \\
\end{tabularx}
\end{projectcard}

\paragraph{Summary.}
In reinforcement learning an agent learns from rewards rather than labels. This project replaces the
policy network of a REINFORCE agent with a parameterised quantum circuit and trains it to balance the
CartPole. It reproduces the design choices that make quantum policies work, input scaling and data
re-uploading, and compares sample efficiency with a classical policy of similar size.

% ===========================================================================
\section{Background and key concepts}
% ===========================================================================
\subsection{Policy gradients}
An agent in state $s_t$ takes action $a_t\sim\pi_\theta(a\,|\,s_t)$ and receives reward $r_t$. REINFORCE
\citep{williams1992} maximises the expected return with the gradient
\begin{equation}
  \nabla_\theta J(\theta)=\E_{\pi_\theta}\Big[\sum_t\nabla_\theta\log\pi_\theta(a_t\,|\,s_t)\,\big(G_t-b\big)\Big],
  \qquad G_t=\sum_{k\ge0}\gamma^kr_{t+k},
\end{equation}
where $b$ is a baseline that reduces variance \citep{sutton2018}.

\subsection{Quantum policies}
\citet{jerbi2021} define a softmax quantum policy
\begin{equation}
  \pi_\theta(a\,|\,s)=\frac{e^{\beta\langle O_a\rangle_{s,\theta}}}{\sum_{a'}e^{\beta\langle O_{a'}\rangle_{s,\theta}}},
\end{equation}
where $\langle O_a\rangle_{s,\theta}$ is the expectation value of an observable assigned to action $a$
after a circuit that encodes $s$ with trainable scaling, $R_X(\lambda_is_i)$, and re-uploads it in
every layer; $\beta$ is a trainable inverse temperature. They found input scaling and re-uploading
essential for performance. \citet{skolik2022} built the corresponding quantum deep Q-learning agent.

\subsection{CartPole}
The cart-pole balancing task \citep{barto1983}, in its Gymnasium version \texttt{CartPole-v1}
\citep{towers2024}, has a 4-dimensional state (cart position and velocity, pole angle and angular
velocity), two actions and a reward of 1 per time step, up to 500 steps. It counts as solved at an
average return of 475 over 100 episodes.

% ===========================================================================
\section{Research questions}
% ===========================================================================
The project is open-ended: these questions are the starting point, not the finish line. Add new
ones to the README and the journal as they come up.

\begin{rqlist}
  \item How many episodes does the quantum policy need to solve CartPole compared with a classical policy of similar size?
  \item Which matters more for the quantum policy: trainable input scaling or data re-uploading?
  \item How sensitive is learning to the number of layers?
\end{rqlist}

% ===========================================================================
\section{Method}
% ===========================================================================
\subsection{Data}
Episodes generated by the agent in \texttt{CartPole-v1}. States are normalised with fixed bounds
before encoding. Each update uses a batch of 10 episodes.

\subsection{Models}
4 qubits, $L$ layers; each layer encodes the scaled state with $R_X(\lambda_is_i)$, applies trainable
$R_Y$/$R_Z$ rotations and a ring of CZ gates. Two-qubit $Z$ observables, one per action, and a
trainable $\beta$ form the softmax policy. Separate Adam learning rates for scaling parameters,
rotation angles and $\beta$; discount $\gamma=0.99$; return normalisation as baseline.

\subsection{Baselines}
A classical policy MLP $4\to h\to2$ with matched parameter count; a random policy; the published
learning curves of \citet{jerbi2021} as a reference.

\subsection{Metrics}
\begin{itemize}
  \item Average return over the last 20 episodes as a function of episodes.
  \item Episodes needed to reach an average return of 475.
  \item Spread across 5 seeds; number of parameters; simulation time.
\end{itemize}

% ===========================================================================
\section{Planned experiments}
% ===========================================================================
Each experiment becomes one \texttt{[[experiments]]} entry in the workspace, so its runs, parameters
and metrics are stored and compared in the portal.

\begin{table}[h]
\centering\small
\begin{tabularx}{\linewidth}{@{}l>{\raggedright}X>{\raggedright}X>{\raggedright\arraybackslash}X@{}}
\toprule
ID & Experiment & What varies & Main metrics \\
\midrule
  E1 & Full quantum policy & $L=5$, scaling + re-uploading & episodes to solve \\
  E2 & No input scaling & $\lambda_i=1$ fixed & learning curve \\
  E3 & No re-uploading & encode once & learning curve \\
  E4 & Depth & $L\in\{2,3,5,7\}$ & episodes to solve \\
  E5 & Q-learning variant & DQN-style agent & sample efficiency \\
\bottomrule
\end{tabularx}
\caption{Experiment plan. Results go into the journal and the results log (\cref{sec:log}).}
\end{table}

% ===========================================================================
\section{Running it in the QML Lab}
% ===========================================================================
The workspace lives in \path{qml_projects/quantum-rl-cartpole/}. Scripts import three helpers from
\path{qml_lab/plots.py}:
\begin{itemize}
  \item \texttt{params()} reads the parameters from the portal's run form (defaults when run by hand);
  \item \texttt{save()} stores figures with the run;
  \item \texttt{log\_metrics()} records numbers that appear in the run table.
\end{itemize}

\begin{lstlisting}
qml_projects/quantum-rl-cartpole/
  project.toml          metadata + experiments
  README.md             goal, open questions, ideas
  experiments/          runnable scripts
  docs/main.tex         this document
\end{lstlisting}

First experiment: This experiment is \emph{planned}: write the script, add the block below to \texttt{project.toml}, then run \texttt{python manage.py sync\_projects}.

\begin{lstlisting}[language={}]
[[experiments]]
slug = "reinforce"
title = "REINFORCE with a softmax quantum policy"
script = "experiments/e01_reinforce.py"
figure = "learning_curve.png"

[experiments.params]
n_layers = 5
episodes = 500
batch_episodes = 10
lr = 0.01
gamma = 0.99
seed = 0
\end{lstlisting}

\begin{lstlisting}[language=bash]
# from the portal folder: pick up changes to project.toml
python manage.py sync_projects
cd qml_projects/quantum-rl-cartpole
# run by hand with the default parameters
PYTHONPATH=../../qml_lab python experiments/e01_reinforce.py
\end{lstlisting}

% ===========================================================================
\section{Quality bar}
% ===========================================================================
\begin{itemize}
  \item A classical baseline with a similar number of trainable parameters, tuned with the same care \citep{bowles2024}.
  \item At least five random seeds per configuration, reported as mean $\pm$ standard deviation.
  \item How results change with the number of qubits.
  \item At least one run under a realistic noise model.
  \item Report full learning curves with seed spread, not only the best run.
\end{itemize}

% ===========================================================================
\section{Open directions}
% ===========================================================================
\begin{itemize}
  \item Other environments (Acrobot, FrozenLake).
  \item Policies estimated from finite shots, as on real hardware.
  \item Run the trained policy on hardware for inference only.
\end{itemize}

% ===========================================================================
\section{Risks and limitations}
% ===========================================================================
\begin{itemize}
  \item RL results have high variance; five seeds is the minimum.
  \item Simulation is slow: every environment step needs a circuit evaluation.
  \item Hyperparameters (learning rates, $\beta$ initialisation) strongly affect results.
\end{itemize}

% ===========================================================================
\section{Results log}
\label{sec:log}
% ===========================================================================
Summarise findings here as they become stable; day-to-day notes belong in the portal journal.

\begin{table}[h]
\centering\small
\begin{tabularx}{\linewidth}{@{}l l >{\raggedright}X >{\raggedright\arraybackslash}X@{}}
\toprule
Date & Run & Finding & Next step \\
\midrule
 & & & \\
 & & & \\
 & & & \\
\bottomrule
\end{tabularx}
\end{table}

\bibliographystyle{harvardqml}
\bibliography{references}

\end{document}
